\subsection{NN-G036: the exact unique 2744-pair optimum for the symmetric closed table}
\paragraph{Statement.}
The minimum cancelling-pair score in the specified symmetric
closed-table presentation is exactly
\[
                                2744.                \tag{G036.1}
\]
The actual polynomial parameter realization attaining it is
unique within that presentation and is NN-G033.
All coefficient degrees and contents, all grading support
locations, and all rational ratios and auxiliary parameters
are allowed, subject to integrality and effectivity of the
actual structure constants. This is not a numerical lower
bound for unrestricted extra-position multiplication tables.

The proof retains its checked stages in their original order.
The first stage forces every competitive $g_1=64$ table into
NN-G034, which has the unique $2744$ attainer. The remaining
ordinary-dimension rows are then eliminated by uniform
integral polynomial arguments. Intermediate intervals and
row counts below describe successive stages, not the final
gate. The last proof and complete projection leave no other
competitive row.

\paragraph{Uniform irreducibility of the small gluing characters.}
Retain $T_j=q^j+q^{-j}$ as an integral polynomial in $D$.
Every mass-eight character in (G034.5) is irreducible
over $\mathbb Q[D]$.
Every mass-ten character there is irreducible except
\[
 W=D+2+2T_j\quad(j\equiv2\pmod4),
\]
in which case
\[
               W=(D-2)B,\qquad B\text{ primitive and irreducible},
               \qquad B\equiv1\pmod2,\quad B(4)=5.   \tag{G036.2}
\]
These statements include all support indices $j$.

Here is an elementary proof. The degree-one choices
are immediate. Every remaining $W$ is primitive,
has leading coefficient two, and reduces to $D$
modulo two. In any nontrivial primitive integral
factorization, one factor has leading coefficient
$\pm1$ and the other $\pm2$. The first factor's
degree is unchanged upon reduction modulo two;
since the product of the reductions is $D$, it
must have degree one. Up to a sign it is $D-r$,
with $r$ even. The other factor reduces to one,
so its constant term is odd.

For mass eight, $W(0)=\pm4$, hence $r=\pm4$.
But $W(4)=8$, and $W(-4)=-4+2T_j(-4)\ne0$.
The last assertion follows from the recurrence
used in NN-G034: at $D=-4$, the absolute values
of $T_j$ are at least six and strictly increase.
There is therefore no such factor.

For mass ten, $W(0)$ is $6$ or $-2$, hence
$r\in\{2,-2,6,-6\}$.
At $D=-2$ the value of $W$ is $2T_j(-2)\ne0$;
at $D=6$ it is positive; and at $D=-6$ the
same recurrence, with initial absolute value eight,
makes $-4+2T_j(-6)$ nonzero.
At $D=2$ the value is $4+2T_j(2)$, which is
zero exactly for $j\equiv2\pmod4$.
This proves that $D-2$ is the only possible factor.
After division by it, $B$ is primitive with leading
coefficient two and reduction one modulo two.
If $B$ factored nontrivially, a factor with odd
leading coefficient would retain positive degree
modulo two, impossible when the product reduces
to one. Thus $B$ is irreducible.
Finally $B(4)=W(4)/2=5$.

\paragraph{The consequence of $g_1=64$.}
The complete mass list NN-G035 gives
$d_1=4,u_1=4$ and $w=8$ or ten in this branch.
Therefore
\[
                 \delta_1=\psi:=D(D-3),\qquad U_1=D.
\]
The primitive $W$ is coprime to $D$ in both small-mass
families. The integral divisibility argument of
NN-G031 gives polynomials $C,Q\in\mathbb Z[D]$ with
\[
 \begin{aligned}
 G_2&=DWC,&G_1&=D\delta_1+D^2C=D^2Q,\\
 C&\equiv D-1\pmod2,&Q&=C+D-3\equiv0\pmod2,\\
 C(4)&=3,&Q(4)&=4.
 \end{aligned}                                                       \tag{G036.3}
\]
No effectivity of $C$ or $Q$ is asserted here.

If $W$ is irreducible, it cannot divide $Q$:
primitivity would make the quotient integral,
so $W(4)=8$ or ten would divide $Q(4)=4$.
Thus $W$ is coprime to $G_1$ and the second-cell
identity implies $U_2=WR$ integrally, with
$R\equiv1\pmod2$.
For $w=8$, the scalar relation
$u_2=12-d_2/2$ gives $4\le u_2\le10$.
Its positive odd ratio $u_2/8$ is therefore one.
The equal-mass lemma of NN-G034 implies $U_2=W$
and then $d_2=8$.
For $w=10$, evenness and positivity of
$u_2=5(30-d_2)/8$ give $d_2=14,u_2=10$.
The same equal-mass lemma gives $U_2=W$.

In the exceptional mass-ten case (G036.2), the
irreducible $B$ cannot divide $Q$, since $B(4)=5$
does not divide four. It is also coprime to $D$.
Thus $U_2=BR$ integrally; modulo two, $R\equiv D$.
The scalar relation still gives $u_2=10$.
Since $\deg B=j-1$ and its leading coefficient is
two, the same leading-coefficient argument as in
NN-G034 forces $\deg R=1$: $R$ cannot be constant
by its parity, and the mass-ten $U_2$ then has
leading coefficient two, so $R$ has leading
coefficient one, which its parity allows only
in degree one.
Therefore $U_2$ has degree $j$ and is the unique
mass-ten character of that degree, namely $W$.
We have proved $U_2=W$ in every case, with no
generic coprimality assumption left unexamined.

\paragraph{Small mass forces the polynomial $C$ itself.}
The two first identities now imply
\[
 \delta_1+\delta_2=(W-D)C.
\]
Put $F=(\delta_1+\delta_2)/2$.
It is an effective symmetric integral polynomial,
dominates $\psi$ coefficientwise, and has mass
six for $w=8$ or nine for $w=10$.
The character $C$ has mass three and satisfies
$C\equiv D-1\pmod2$.

If $W=D+4$ or $D+6$, then $F=2C$ or $3C$.
In particular $C$ is effective. Since $F\ge\psi$,
its coefficients at $q^{\pm1},q^{\pm2}$ force
those of $C$ to be at least one. This would give
$C(4)\ge4$, a contradiction.

For the remaining mass-eight characters,
$F=T_jC$. Its leading coefficient is positive.
If $\deg C>1$, that coefficient in $C$ is even
and at least two by parity. The degree of $F$
then exceeds two, so its outer pair contributes
at least four mass units disjointly from the four
units already required by $\psi$. This contradicts
$F(4)=6$. Thus $\deg C=1$.

For mass ten, $F=(1+T_j)C$. If $\deg C>1$,
the same argument produces an outer pair of
coefficient at least two outside $\psi$.
Since $F(4)=9$, the only possibility is exactly
\[
                         F=\psi+1+2T_k
\]
for some $k>2$: four mass units come from $\psi$,
four from that outer pair, and the remaining
unit must be central. This is impossible modulo
two at $D=1$. Its displayed value is odd there,
whereas $F=(1+T_j)C$ is even because $C(1)$ is
even. Therefore $\deg C=1$ also in this case.

Write $C=aD+b$, with $a$ positive odd and
$4a+b=3$. Its symmetric central coefficient is
$3-2a$ and its $q^1$ coefficient is $a$.
If $w=8,j=1$, the $q^1$ coefficient of $T_1C$
is $3-2a$ and must be at least one, forcing $a=1$.
If $w=10,j=1$, the corresponding coefficient of
$(1+T_1)C$ is $3-a$, again forcing the positive
odd $a$ to be one.
For $j\ge2$, the leading pair is outside $\psi$;
an $a\ge3$ would give total mass at least
$4+2a\ge10$, exceeding six or nine.
Thus in every case
\[
                  C=D-1,\qquad Q=2D-4,\qquad
                  G_1=2D^2(D-2).                    \tag{G036.4}
\]
The three fixed first-cell data of NN-G034 are now
forced, rather than assumed. That theorem proves
that the sole competitive realization in this
branch is NN-G033, of cost exactly $2744$.

\paragraph{The improved finite gate and lower bound.}
Remove from the NN-G035 certificate the five base
rows with $g_1=64$. Every strict improvement must
lie in the remaining $38$ rows, with $348$
auxiliary increment pairs and $20$ first-cell
mass triples. The least remaining relaxed score
is $1932$; the exact projection checker verifies
these four numbers. The removed branch has the
unique actual score $2744$ just proved. This establishes
the intermediate lower bound $1932$; the combined
refinement below strengthens it to $2572$.
\hfill$\square$

\paragraph{Geometric translation and scope.}
Within the NN-G010 realization by graded sign-representation
kernels, a competitive first-cell bottom-action object
of ordinary dimension $64$ determines its entire character
and, in turn, the unique $2744$ completion in this closed
presentation. An improvement requires dimension at least
$72$ for that object and initially obeys $20$
dimension triples, reduced to $15$ by the addendum below. This statement has no bound on the
grading support of the other surviving objects.
The irreducibility argument is a statement in the
integral grading-character ring, not a claim of
irreducibility of the representations themselves.

The dictionary is SGA 3, I, 4.7.2--4.7.3; the character
evaluation at $q=-1$ and the parity reductions are
integral restrictions of this ring. Actual composition
uses archive cote 2, scan 17 through NN-G010.
The proof neither identifies a canonical GCT basis nor
applies this lower bound to the general extra-position
problem. At this intermediate stage only a lower bound was established.
The later exclusions below close the entire closed-table gate;
the extra-position optimization remains separate.


\paragraph{Combined refinement: the 2572 lower bound for closed tables.}
This independently reviewed addendum strengthens the
preceding intermediate bound. The full argument below
excludes twelve more necessary base rows and proves the
closed-table interval $[2572,2744]$.
Its first step excludes the two
NN-G035 base rows
\[
 (w,g_1,g_2,d_1,d_2,u_1,u_2)
   =(8,72,96,6,6,4,8),\quad(8,80,96,8,4,4,8),
\]
with every allowed auxiliary increment, not only their
least auxiliary masses. The mass-ten rows with first-cell
triple $(80,8,4)$ are not excluded by this argument.

In either row, $U_1=D$ and the earlier integral quotient
argument gives $G_1=DH$, $G_2=DWC$,
$H=\delta_1+DC$, $C\equiv D-1\pmod2$, and $C(4)=3$.
The mass-eight $W$ is irreducible and coprime to $D$.
It cannot divide $H$, whose mass is respectively $18$
or $20$, because eight divides neither number.
Thus $W$ is coprime to $G_1$ and both second-cell
identities give $U_2=WR,V_2=WZ$ integrally.
Their parities give $R\equiv Z\equiv1$.
The equal-mass lemma forces $U_2=W$, and hence
\[
             \delta_1+\delta_2=(W-D)C.               \tag{G036.R1}
\]

For the first row, both cells have their unique
mass-six character $\delta_i=\psi+2$, where
$\psi=D(D-3)$. If $W=D+4$, (G036.R1) would give
$C=(\psi+2)/2$, which is not an integral polynomial.
Otherwise $W=D+2T_j$ and
$T_jC=\psi+2=(D-1)(D-2)$.
Since $C\equiv D-1$ has positive degree, this forces
$j=1$, $C=D-1$, $W=3D-4$ and $H=2(D-1)^2$.
Now
\[
                 A_2=D^2WC-DHZ.
\]
If $Z$ has degree one, its leading coefficient is
positive even, so the leading coefficient of $A_2$
is at most $3-4<0$. Higher degrees also fail.
Thus $Z$ is constant, but then $V_2=WZ$ has degree
one and cannot dominate its required $q^2$ coefficient
two. This excludes the first row.

For the second row, $\delta_2=\psi$ and
$F=(\delta_1+\delta_2)/2$ is effective, has mass six,
and dominates $\psi$. The mass-eight part of the
small-mass argument in NN-G036 forces $C=D-1$;
the only possible $W$ are $3D-4$ and $D+2T_2$.
For $W=3D-4$ one has $\delta_1=\psi+4$ and
$H=2D(D-2)+4$. Its degree and leading coefficient
give the same contradiction for $Z$ as in the first row.

For $W=D+2T_2$, the forced polynomials are
\[
 \delta_1=\psi+2T_3,\qquad
 H=2D(D-2)+2T_3,\qquad C=D-1.
\]
Here $D^2WC$ has degree five and leading coefficient
two, while $DH$ has degree four and leading coefficient
two. A nonconstant $Z\equiv1\pmod2$ would again give
a negative leading coefficient in $A_2$.
Consequently $Z$ is a positive odd integer.
The scalar relations give $v_2=8Z$ and
$a_2=384-10v_2$. Since $v_2\ge16$ and $a_2\ge16$,
the only possible odd value is $Z=3$, so
$v_2=24,a_2=144$.

The first auxiliary relation is $a_1=320-12v_1$.
It gives $v_1\le26$; and substitution of every other
now-fixed mass into the total gives
\[
                    T=9760+148(26-v_1).
\]
Since $v_1$ is even, the cap $T\le9912$ forces
$v_1=26$ and $a_1=8$.
We show that these masses cannot satisfy coefficientwise
domination.

Effectivity of $A_1=D^2H-D(D-1)V_1$ gives
$\deg V_1\le3$. Write $V_1=(d,a,b,c)$ in symmetric
Laurent notation. Domination and parity require
$d\ge6$ even, $a\ge5$ odd, $b\ge0$ even, and
$c\ge0$ even. The leading coefficient of $A_1$
also gives $c\le2$.
The mass condition is $d+2a+2b+2c=26$.
The relevant coefficients of $A_1$, after eliminating
$d$, are
\[
 \begin{array}{c|l}
 q^0&-44+2a+6b+8c,\\
 q^1&-24+a+3b+5c,\\
 q^2&14-a-2b-c,\\
 q^3&24-a-3b-4c.
 \end{array}
\]
The target lower bounds at degrees zero and three,
respectively two and one, force
$a+3b+4c=23$. The degree-one target then gives
$c-1\ge1$, so $c=2$ and $a+3b=15$.
The degree-two target gives $b-3\ge1$, hence $b\ge4$.
This implies $a\le3$, contradicting $a\ge5$.
The second row is therefore impossible as well.
\hfill$\square$

Deleting these two rows alone leaves $36$ base rows,
$311$ auxiliary increment pairs, and $19$ first-cell
mass triples, with least relaxed score $1956$.
The following uniform arguments exclude three more
mass blocks.

\paragraph{A common-factor observation.}
The identities used throughout are
\[
 WS=G_1G_2,\qquad U_1G_2=W(G_1-D\delta_1),\qquad
 U_2G_1=W(G_2-D\delta_2).                            \tag{G036.R2}
\]
If a primitive irreducible factor $P$ of $W$ divides
neither $U_1$ nor $U_2$, then it divides both $G_1$
and $G_2$. Indeed it divides at least one from the
first identity, and the appropriate one of the next
two forces it into the other. All divisibilities
first obtained over $\mathbb Q[D]$ are integral by
primitivity. In particular divisibility by $P$ forces
$P(4)$ to divide the corresponding ordinary mass.

\paragraph{The entire mass block $(w,g_1,g_2)=(10,80,80)$.}
Here $u_i=10-d_i/2$, so $u_i\in\{4,6,8\}$.
If $W$ is irreducible, its mass ten divides neither
$u_i$, and (G036.R2) gives $W\mid G_1,G_2$.
If $W=(D-2)B$ is the exceptional case of (G036.2),
the primitive irreducible $B$ has mass five, again
dividing neither $u_i$. Thus write $G_i=BP_i$.
The second and third identities give
\[
 U_1P_2=(D-2)(BP_1-D\delta_1),\qquad
 U_2P_1=(D-2)(BP_2-D\delta_2).
\]
An effective character of mass four, six or eight,
dominating and congruent to $D$, has nonzero value
at $D=2$: these values are $2$, $4$, or one of
$6,2,-2$, respectively. Therefore $D-2$ divides
both $P_i$. We have proved $W\mid G_1,G_2$ also
in this exceptional case.

Now write $G_i=WP_i$ in every case. Then
\[
 P_1(4)=P_2(4)=8,\qquad P_1\equiv0,\quad
 P_2\equiv D(D-1)\pmod2.
\]
The polynomial $P_2$ cannot be divisible by $D$:
its integral quotient would have value two at $D=4$
but would reduce to $D-1$ modulo two, whose value
there is odd. Thus $P_2(0)\ne0$.
The divided identities at $D=0$ show also
$P_1(0)\ne0$, because $W(0)\ne0$.
Multiplying them gives
\[
                        W(0)^2=U_1(0)U_2(0).        \tag{G036.R3}
\]
For the mass-ten characters, $W(0)$ is six or
minus two. For the possible $U_i$, the values are
zero, two, four or minus four. Equality (G036.R3)
is possible only when $W(0)=-2$ and both $U_i(0)=2$.
Consequently
\[
 U_1=U_2=D+2,\qquad d_1=d_2=8,\qquad
 W=D+2+2T_j\quad\hbox{with }j\hbox{ odd}.             \tag{G036.R4}
\]
Each cell character is then $\delta_i=\psi+2R_i$,
where $R_i$ has mass two and is either the constant
two or some $T_k$.

Subtracting the two divided identities in (G036.R2),
and putting $E=R_1-R_2$, gives
\[
 A_0(P_1-P_2)=E,\qquad
 A_0=\frac{D+2+T_j}{D}.                             \tag{G036.R5}
\]
The quotient is integral since $j$ is odd and
$T_j(0)=-2$. The difference $P_1-P_2$ is nonzero,
as its reduction is $D(D-1)$, so $E\ne0$.
Every nonzero difference between two members of
$\{2,T_1,T_2,\ldots\}$ has leading coefficient
$1$ or $-1$, and its value at $D=0$ has absolute
value at most four.

If $j=1$, $A_0=2$, contradicting this odd leading
coefficient of $E$.
For odd $j\ge3$, $A_0$ is monic of degree $j-1$.
Thus $P_1-P_2$ has leading coefficient $\pm1$;
its reduction $D(D-1)$ forces its actual degree
to be exactly two. Write it $aD^2+bD+c$, with
$a=\pm1$, $b$ odd and $c$ even.
The recurrence $T_{j+1}=(D-2)T_j-T_{j-1}$ gives
\[
 T_j'(0)=(-1)^{j-1}j^2,\qquad A_0(0)=1+j^2.
\]
For completeness, the derivative recurrence is
$r_{j+1}=2(-1)^j-2r_j-r_{j-1}$ with
$r_0=0,r_1=1$; the displayed formula satisfies it
by direct substitution.
Evaluating (G036.R5) at zero now gives
$|(1+j^2)c|\le4$, so $c=0$.
But $P_1(4)=P_2(4)$ then implies
$16a+4b=0$, making $b=-4a$ even, a contradiction.
This excludes the entire mass block, including
all six base rows and all their auxiliary choices.

\paragraph{The entire mass block $(w,g_1,g_2)=(8,96,96)$.}
Here $u_i=8-d_i/3$ is four or six.
The mass-eight $W$ is irreducible, and its mass
divides neither $u_i$, so (G036.R2) again gives
$G_i=WP_i$. Now $P_i(4)=12$, with
$P_1\equiv0$ and $P_2\equiv D(D-1)$ modulo two.
The polynomial $P_1$ cannot be divisible by $D$:
the quotient would have even coefficients but
ordinary mass $12/4=3$.
The divided identities at zero consequently
force both $P_i(0)\ne0$ and give (G036.R3).
But $W(0)^2=16$, whereas $U_i(0)$ are zero or two,
whose products are zero or four. This contradiction
excludes both base rows in this block.

\paragraph{The two asymmetric mass-ten blocks.}
The finite certificate has one base row for
$(w,g_1,g_2)=(10,76,100)$, with
$(d_1,d_2,u_1,u_2)=(4,6,6,10)$, and its index-swapped
row with $(g_1,g_2)=(100,76)$.
Let $G_-$ denote the smaller action character and
$G_+$ the larger. Their cell characters are
respectively $\psi$ and $\psi+2$, and the smaller
index has $U_-=D+2$.

If $W$ is irreducible, it cannot divide $G_-$,
since ten does not divide $76$.
It therefore divides $G_+$ and then $U_+$ by
(G036.R2). Since $u_+=10$, the equal-mass lemma
forces $U_+=W$.
If $W=(D-2)B$, the irreducible mass-five $B$ cannot
divide $G_-$, so the same argument gives $B\mid U_+$.
Writing $U_+=BR$, its reduction is $R\equiv D$.
The leading-coefficient argument from NN-G036 forces
$\deg R=1$ and then the mass-ten classification gives
$U_+=W$ in this case as well.

The larger-index identity now gives
$G_+=G_-+D(\psi+2)$.
Substitution into the smaller-index identity yields
\[
 (W-D-2)G_-=D\bigl(W\psi+(D+2)(\psi+2)\bigr).       \tag{G036.R6}
\]
If $W=D+6$, its denominator is four and the resulting
$G_-$ has leading coefficient $1/2$, contradicting
integrality.
For $W=D+2+2T_j$, equation (G036.R6) gives
\[
 G_-=D\psi+\frac{N(D)}{T_j(D)},\qquad
 N(D)=D(D+2)(D^2-3D+1).
\]
Integrality forces the monic $T_j$ to divide $N$.
Thus $j\le4$. None does:
$N(2)=-8$ excludes $T_1=D-2$ and also $T_3$,
which contains $T_1$;
the remainder modulo $T_2=D^2-4D+2$ is $16D-10$;
and $T_4=D^4-8D^3+20D^2-16D+2$ is not the
monic degree-four polynomial $N$.
This excludes both asymmetric blocks, uniformly
over all original support indices.
\hfill$\square$

\paragraph{The resulting finite dimension gate.}
In addition to the five rows already removed in
NN-G036, this addendum removes two isolated mass-eight
rows, six rows in the $(10,80,80)$ block, two in the
$(8,96,96)$ block, and the two asymmetric rows.
The remaining complete necessary projection has
$26$ base rows, $42$ auxiliary increment pairs,
and $15$ first-cell mass triples. Its least relaxed
score is $2572$. The proved exact attainer still
has score $2744$, so the closed-table minimum lies
in $[2572,2744]$.

In the NN-G010 geometry, these are exclusions of
ordinary-dimension patterns for actual composing
graded sign-representation objects. The tests at
$D=0$ are integral character evaluations at $q=-1$;
the derivative there is a character-polynomial
calculation, not a Weil-weight assertion.
They preserve the SGA 3, I, 4.7.2--4.7.3 and
archive cote 2, scan 17 dictionary of NN-G036.
The surviving rows are not declared realizable,
and this lower bound is not transferred to arbitrary
extra-position lifts or canonical GCT multiplicities.


% Further NN-G036 refinement: root whole-proof PASS 2026-11-22T00:03:47Z.
\paragraph{The last mass-eight competitor.}
The remaining mass-eight base row has
\[
 (g_1,g_2,d_1,d_2,u_1,u_2)=(128,80,12,4,8,4).
\]
Thus $\delta_2=\psi$ and $U_2=D$. The integral divisibility
argument already used in NN-G035 gives
\[
 G_1=DWC,\qquad G_2=D(\psi+DC),\qquad
 C\equiv0\pmod2,\quad C(4)=4.
\]
For clarity, $DG_1=W(G_2-D\psi)$ and $\gcd(W,D)=1$
first give $G_1=WQ$ and $G_2=D(Q+\psi)$ integrally.
The other cell identity at zero forces $G_1(0)=0$,
so $Q(0)=0$ and $Q=DC$. Reducing $G_1=DWC$ modulo
two gives the displayed parity of $C$.
The irreducible mass-eight $W$ cannot divide $G_2$:
it is coprime to $D$, whereas the integral polynomial
$\psi+DC$ has value twenty at four, not divisible by eight.
Consequently $U_1G_2=W(G_1-D\delta_1)$ forces $W\mid U_1$.
The equal-mass lemma gives $U_1=W$. Cancelling now yields
\[
 \delta_1+\delta_2=(W-D)C.
\]
Write $C=2H$ with $H\in\mathbb Z[D]$ and $H(4)=2$.
The effective character
$F=(\delta_1+\delta_2)/2=(W-D)H$ has mass eight and
dominates $\psi$.
If $W=D+4$, then $F=4H$, so $H$ is effective;
domination forces its coefficients at both positive
exponents one and two to be at least one. Its mass
would therefore be at least four, contradicting two.
Otherwise $W=D+2T_j$ and $F=2T_jH$ has all Laurent
coefficients even. Its coefficients at exponents
$1,-1,2,-2$ must all be at least two. These exhaust
the mass eight, hence $F=2\psi$ and $T_jH=\psi$.
Since $T_j$ is monic of degree $j$, this requires
$j\le2$. But $T_1=D-2$ does not divide $\psi$, and
$T_2=D^2-4D+2$ is not the monic quadratic $\psi$.
This excludes the last mass-eight competitor.
Together with the preceding exclusions and NN-G034,
every competitive mass-eight closed table is the unique
$2744$ attainer. No bound on its initial support indices
was imposed.

\paragraph{The mass-ten block with equal action masses one hundred.}
The row is
\[
 (g_1,g_2,d_1,d_2,u_1,u_2)=(100,100,10,10,6,6).
\]
Thus both $U_i=D+2$. If $W$ is irreducible, its mass
ten divides neither $u_i$, so (G036.R2) forces
$W\mid G_1,G_2$. In the exceptional factorization
$W=(D-2)B$, the irreducible $B$ has mass five and
again divides neither $U_i$; it divides both $G_i$.
After cancelling $B$ in the cell identities, evaluation
at $D=2$, where $U_i(2)=4\ne0$, forces the remaining
factor $D-2$ into both $G_i$. Thus in every case
\[
 G_i=WP_i,\qquad P_i(4)=10,\qquad
 P_1\equiv0,\quad P_2\equiv D(D-1)\pmod2.
\]
The polynomial $P_2$ cannot be divisible by $D$,
since its quotient would have nonintegral value $10/4$
at four. The divided cell identities at zero then
force both $P_i(0)\ne0$ and give $W(0)^2=4$.
Consequently $W=D+2+2T_j$ with $j$ odd.
Write $\delta_i=\psi+2R_i$. Now $R_i$ is an effective
character of mass three, hence is either three or
$1+T_k$. A nonzero difference $E=R_1-R_2$ still
has leading coefficient $\pm1$ and $|E(0)|\le4$.
It is nonzero because the parity difference of
$P_1,P_2$ is nonzero.
Subtracting the divided identities gives exactly
\[
 \frac{D+2+T_j}{D}(P_1-P_2)=E.
\]
For $j=1$, the factor is two, contradicting the odd
leading coefficient of $E$. For odd $j\ge3$, it is
monic, with value $1+j^2$ at zero. Hence
$P_1-P_2=aD^2+bD+c$, where $a=\pm1$, $b$ is odd,
and $c$ is even: the degree follows from its leading
coefficient and its reduction $D(D-1)$.
The bound on $E(0)$ forces $c=0$.
Equal values of the two $P_i$ at four then give
$16a+4b=0$, contradicting that $b$ is odd.
The row is impossible.

\paragraph{The two mass-ten blocks with action masses eighty-four and one hundred.}
In either orientation the smaller action has
$\delta_-=\psi+2$, $U_-=D+2$, while the larger
has $\delta_+=\psi$ and $u_+=10$.
An irreducible $W$ of mass ten cannot divide $G_-$
of mass eighty-four; (G036.R2) consequently forces
$W\mid G_+$ and $W\mid U_+$, whence $U_+=W$
by the equal-mass lemma.
If $W=(D-2)B$, the mass-five irreducible factor
$B$ likewise cannot divide $G_-$. Thus $B\mid U_+$.
Writing $U_+=BR$, reduction gives $R\equiv D\pmod2$.
The leading coefficient argument in NN-G036 forces
$\deg R=1$; the mass-ten classification then gives
$U_+=W$ also in this case.
The larger-index identity now says $G_+=G_-+D\psi$.
Substituting in the smaller-index identity gives
\[
 (W-D-2)G_-=D\bigl(W(\psi+2)+(D+2)\psi\bigr).
\]
For $W=D+6$ this would make the leading coefficient
of $G_-$ equal to one half. Otherwise
$W=D+2+2T_j$ and
\[
 G_-=D(\psi+2)+\frac{N(D)}{T_j(D)},\qquad
 N=D(D+2)(D^2-3D+1).
\]
The same integral divisibility obstruction as in
(G036.R6) applies: $T_j\mid N$ requires $j\le4$,
but $N(2)=-8$, $T_1\mid T_3$,
$N\bmod T_2=16D-10$, and $N\ne T_4$.
Both orientations are excluded, with every auxiliary
increment allowed by their scalar relaxation.
\hfill$\square$

\paragraph{Scope and geometric content.}
These four base-row exclusions concern the ordinary
dimensions of actual graded $C_2$-representation
structure objects in the matrix-kernel category of
NN-G010, with the same symmetric closed multiplication
table and character constraints. In particular, within
the mass-eight gluing stratum the actual smallest
effective lift is the $2744$-pair object of NN-G033,
uniquely within that presentation. The integral torus
character interpretation is SGA~3, I, 4.7.2--4.7.3;
kernel composition is the supplied archive transcription,
cote~2, scan~17. No new claim about a canonical GCT
representation family or a different multiplication
table follows.

\paragraph{A mass-twelve gluing lemma.}
Throughout, $D=q+2+q^{-1}$ and $T_j=q^j+q^{-j}$.
Every effective $W\ge D$ with $W\equiv D\pmod2$ and $W(4)=12$
is one of
\[
 D+8,\quad D+4+2T_j,\quad D+4T_j,\quad
 D+2T_i+2T_j\ (i<j).
\]
Indeed $(W-D)/2$ is an effective symmetric character of mass four,
and these are its complete partitions into a central coefficient
and paired nonzero exponents. Hence $W(0)\in\{-8,0,8\}$.
If $W(0)=0$, write $W=DH$. Then
\[
 H\in\mathbb Z[D],\quad H(4)=3,\quad H\equiv1\pmod2,
 \quad H(0)\ne1,-1.
\]
Either $H=3$ (the case $W=3D$), or $H$ is primitive and
irreducible over $\mathbb Q[D]$, with leading coefficient two.
For the last assertion, every nonconstant case has leading
coefficient two and odd constant coefficient. In a nontrivial
primitive integral factorization, one factor has odd leading
coefficient and would keep its positive degree modulo two,
contradicting that the product reduces to one.
For the assertion at zero, the zero-valued cases in the list are
$D+4+2T_o$ with $o$ odd and $D+2T_o+2T_e$ with $o$ odd and
$e$ even. The derivative identity from (G036.R5) gives
\[
 H(0)=1+2o^2\quad\hbox{or}\quad 1+2(o^2-e^2).
\]
Neither is one or minus one: equality would require equal
squares of opposite parity or $e^2-o^2=1$, impossible for
positive distinct integers $e,o$. This also includes $o=1$
and the constant $H=3$.

\paragraph{The nonzero-value branch.}
For any effective $U\ge D$ with $U\equiv D\pmod2$ and mass $u$,
\[
 |U(0)|\le u-4,
\]
because $(U-D)/2$ is effective and its absolute value at $q=-1$
is at most its mass. In every remaining mass-twelve row of
NN-G036, $(u_1-4)(u_2-4)<64$, as follows directly from the
complete listed mass tuples. If $W(0)=8$ or $-8$, the identities
\[
 WG_1-U_1G_2=DW\delta_1,\qquad
 WG_2-U_2G_1=DW\delta_2
\]
have an invertible coefficient matrix at $D=0$ over $\mathbb Q$.
Thus $D$ divides both $G_i$ in $\mathbb Z[D]$.
Using $G_1\equiv0$ and $G_2\equiv D^2(D-1)\pmod2$, their
quotients by $D$ have even evaluation at $D=4$. Consequently
both $g_i$ are divisible by eight. This excludes any mass row
containing $84$ in this branch.
If also $\delta_1=\delta_2=\psi=D(D-3)$, divide the same
identities by $D$ and apply the invertible matrix at zero once
again. Both $G_i$ are divisible by $D^2$. Since $WS=G_1G_2$
and $W(0)\ne0$, the polynomial $S$ is divisible by $D^4$.
In particular, the mass block $(g_1,g_2)=(80,96)$ or $(96,80)$,
where $s=640$, is impossible: $4^4$ does not divide $640$.

\paragraph{The two blocks with action masses $84,88$.}
Their action masses are $84,88$ in either order, with associated
$u$-masses six and eight. The preceding paragraph leaves only
$W=DH$. If $H$ is nonconstant, it cannot divide the mass-six
character $D+2$ (the integral leading coefficient would have
to divide one), nor a mass-eight character (its value three
would have to divide eight). The common-factor argument of
(G036.R2) forces $H$ to divide both $G_i$, contradicting that
three does not divide $88$.
If $H=3$, reduction modulo three supplies the same argument
in the integral domain $\mathbb F_3[D]$. Both $U_i$ have
nonzero reduction: one is $D+2$, and the other has mass eight.
The three identities (G036.R2) then force both $G_i$ to have
all coefficients divisible by three, again contradicting
mass $88$. Thus both base rows are impossible.

\paragraph{The block with both action masses $84$.}
Here $U_1=U_2=D$ and $d_1=d_2=14$. Again only $W=DH$
remains. The common-factor argument (or reduction modulo
three when $H=3$) gives $G_i=HP_i$ integrally. Their masses
are $P_i(4)=28$, and their reductions are
$P_1\equiv0$, $P_2\equiv D^2(D-1)\pmod2$.
Cancelling $DH$ in the two cell identities gives
\[
 P_2=HP_1-D\delta_1,\qquad
 P_1=HP_2-D\delta_2.
\]
Since $H(0)^2\ne1$, evaluating at zero forces
$P_1(0)=P_2(0)=0$. In particular $P_1/D$ is an integral
polynomial with every coefficient even, but its evaluation
at four is $28/4=7$. This is a contradiction.

\paragraph{The two smallest mass-twelve rows.}
Suppose $(g_1,g_2)$ is $(80,96)$ or $(96,80)$ and both
cells are $\psi$. Let $G_-$ have mass eighty and $G_+$
have mass ninety-six; their $U$-masses are eight and twelve.
The nonzero-value branch was already excluded, so $W=DH$.
For nonconstant $H$, its value three cannot divide eighty,
so $H\nmid G_-$. The identities (G036.R2) force
$H\mid G_+$ and $H\mid U_+$. Every mass-twelve $U_+$
has value at zero equal to zero or $\pm8$. Since $H(0)$ is
odd and is neither one nor minus one, its divisibility
forces $U_+(0)=0$. Write $U_+=DJ$. The classification above
applies to $J$: it is either three or primitive irreducible
of leading coefficient two. Since $H\mid J$, the latter
case is forced and $J=H$. Thus $U_+=W$.
For $H=3$, reduce $U_+G_-=3D(G_+-D\psi)$ modulo three.
Since $G_-$ has mass eighty, its reduction is nonzero, so
all coefficients of $U_+$ are divisible by three. Its odd
coefficients at $q,q^{-1}$ are at least three each, and its
even central coefficient is at least six. These exhaust its
mass twelve, so $U_+=3D=W$ also in this case.

The larger-cell identity now gives $G_+=G_-+D\psi$.
The other identity becomes
\[
             (W-U_-)G_-=D\psi(W+U_-).
\]
The mass-eight classification gives $U_-(0)=4$ or $-4$,
while $W(0)=0$. Hence $D^2\mid G_-$.
If $G_-=G_1$, its quotient by $D^2$ has even coefficients
but mass $80/16=5$, a contradiction.
If $G_-=G_2$, use the necessary mod-four identity
$G_2-D\delta_2\equiv G_1\pmod4$ from NN-G030.
Here $\delta_2=\psi$ and $G_1=G_2+D\psi$, so it would
require $2D\psi\equiv0\pmod4$, which is false.
This excludes both rows, including every auxiliary increment.

These five exclusions use neither a degree bound nor
positivity of any quotient. They retain precisely the
closed-table character dictionary and scope of NN-G036.

\paragraph{Intermediate combined gate at 00:03:47 UTC.}
The additional four small-gluing and five mass-twelve exclusions
leave exactly $17$ necessary base rows, $26$ auxiliary increment
pairs, and $8$ first-cell mass triples. The least remaining
relaxed score is $2604$, so the intermediate closed-table
interval at this stage is $[2604,2744]$. Every strict improvement has $w=10$
or $12$; in the mass-ten case its first data must be
$(g_1,g_2,d_1,u_1)=(80,120,8,4)$, with
$d_2\in\{6,10,14,18\}$. The earlier numerical intervals in
this subsection are successive intermediate bounds.
The exact finite projection is
\texttt{verify\_second\_competitive\_refinement.py};
the preceding polynomial arguments justify every deletion
without a support-degree bound.

\paragraph{Geometric conclusion at the intermediate stage.}
Within the specified closed presentation of composing graded
$C_2$-representation objects from NN-G010, every effective
lift has at least $2604$ cancelling trivial/sign pairs and
one with $2744$ pairs exists. In the gluing-dimension-eight
stratum, that $2744$ object is the unique minimum in this
presentation. A strict improvement can occur only in the
seventeen displayed ordinary-dimension rows, whose actual
character realization remains a further condition.
The bounds concern genuine effective representation objects,
using SGA~3, I, 4.7.2--4.7.3 and the archive kernel-product
transcription at cote~2, scan~17. They are not statements
about other multiplication-table shapes or canonical GCT
multiplicities.


% NN-G036 further uniform refinement; root whole-proof PASS 2026-11-22T00:06:53Z.
\paragraph{The entire remaining mass-ten gluing branch.}
The reviewed finite projection leaves four mass-ten base rows:
\[
 \begin{aligned}
 &g_1=80,\quad g_2=120,\quad d_1=8,\quad U_1=D,\\
 &(d_2,u_2)\in\{(6,12),(10,10),(14,8),(18,6)\}.
 \end{aligned}
\]
Every allowed auxiliary increment in these rows has
\[
 a_1=8,\qquad v_1=26,\qquad v_2\in\{56,58\}.
\]
Equality with the first auxiliary's absolute trace mass
therefore fixes the entire polynomial, not only its mass:
\[
 A_1=D(D-2)(D-3)
   =2+q+q^{-1}+q^2+q^{-2}+q^3+q^{-3}.                 \tag{G036.T1}
\]
We exclude all four rows, with arbitrary support locations.

The integral $U_1=D$ argument gives
\[
 \begin{aligned}
 G_2&=DWC,\qquad G_1=DH,\qquad H=\delta_1+DC,\\
 C&\equiv D-1\pmod2,\qquad H\equiv0\pmod2,\\
 C(4)&=3,\qquad H(4)=20.
 \end{aligned}                                                   \tag{G036.T2}
\]
The mass-eight cell is necessarily
$\delta_1=\psi+2R$, where $R$ is either two or $T_k$.
In particular $\delta_1(0)$ is four or minus four.

First suppose $W$ is irreducible. If it did not divide
$H$, it would be coprime to $G_1=DH$, because $W(0)\ne0$.
The second auxiliary identity
$V_2G_1=W(DG_2-A_2)$ would force $W\mid V_2$ integrally.
Evaluation at four would then give $10\mid v_2$, contrary
to $v_2=56$ or $58$. Thus $W\mid H$.

In the exceptional case $W=(D-2)B$ of (G036.2),
$B$ is primitive irreducible, $B(4)=5$,
$B\equiv1\pmod2$, and $B(0)=-3$.
If $B\nmid H$, the same auxiliary identity gives
$B\mid V_2$, contradicting $5\nmid v_2$.
Therefore $H=2BK$ for an integral polynomial $K$;
the factor two follows by reducing $H$ modulo two.
But then $H(0)$ is divisible by six, whereas
$H(0)=\delta_1(0)=\pm4$. The exceptional case is impossible.

We are left with the irreducible case, where
\[
 H=2WK,\qquad K\in\mathbb Z[D],\qquad K(4)=1.
\]
Here the integral quotient by $W$ is even because
$\overline W=D$ and $\overline H=0$ in the domain
$\mathbb F_2[D]$. If $W(0)=6$, the identity
$H(0)=2W(0)K(0)=\pm4$ is impossible. Hence
\[
 W=D+2+2T_j,\quad j\ {\rm odd},\qquad W(0)=-2,
 \qquad K(0)\in\{1,-1\}.
\]
Every integer polynomial satisfies $K(4)\equiv K(0)\pmod4$.
Since $K(4)=1$, this forces $K(0)=1$.
Consequently $\delta_1(0)=H(0)=-4$, so
$R=T_k$ with $k$ odd. At $D=2$ the Chebyshev
evaluations are therefore
\[
 W(2)=4,\qquad \delta_1(2)=\psi(2)+2T_k(2)=-2.
                                                                  \tag{G036.T3}
\]

The first auxiliary identity, after cancelling $W$
and then $D$, is
\[
 C V_1=2DWK-A_1/D.
\]
By (G036.T1), evaluation at two gives
\[
                  C(2)V_1(2)=16K(2).               \tag{G036.T4}
\]
The integer $C(2)$ is odd by (G036.T2).
Therefore $16$ divides $V_1(2)$.
The coefficientwise lower bound for $V_1$ is the
effective character $V_{\min}=6+5(q+q^{-1})$,
of mass sixteen. Its parity and mass twenty-six imply
$V_1=V_{\min}+2P$ for an effective symmetric integral
character $P$ of mass five. At $D=2$, equivalently
$q=i$, the triangle inequality gives
$|P(i)|\le5$ and hence
\[
                         -4\le V_1(2)\le16.
\]
Its only multiples of sixteen in this interval are
zero and sixteen.
If $V_1(2)=0$, (G036.T4) gives $K(2)=0$.
The monic polynomial $D-2$ would divide $K$ integrally,
so its value two at four would divide $K(4)=1$,
which is impossible.
Thus $V_1(2)=16$, and (G036.T4) gives $K(2)=C(2)$.
But (G036.T2)--(G036.T3) also give
\[
 C(2)=\frac{2W(2)K(2)-\delta_1(2)}2
      =4K(2)+1.
\]
This would require $-3K(2)=1$, contradicting integrality.
Every mass-ten competitor is excluded.
\hfill$\square$

\paragraph{Meaning of this refinement.}
The closed-table optimum is consequently attained at
gluing dimension eight unless a strictly smaller object
occurs in the remaining dimension-twelve stratum.
The proof uses restriction of integral torus characters
at $q=-1$ and $q=i$, together with ordinary dimensions
at $q=1$; it assumes no Weil-weight estimate and no
effectivity of $C,H/W$, or $K$.
Its geometric objects and source dictionary remain the
actual composing graded sign representations of NN-G010,
SGA~3, I, 4.7.2--4.7.3, and the supplied kernel-product
transcription at archive cote~2, scan~17.

\paragraph{Projection after the mass-ten exclusion.}
Deleting its four base rows and six auxiliary increments leaves
$13$ necessary base rows, $20$ auxiliary increments and
$7$ first-cell mass triples. Every strict improvement on
$2744$ now has gluing dimension twelve. The least relaxed
score remains $2604$. The exact row projection is certified
separately by \texttt{verify\_mass10\_final\_exclusion.py}.


% Further NN-G036 refinement; root whole-proof PASS 2026-11-22T00:12:35Z.
\paragraph{A mass-six action gives integral divisibility.}
Suppose $W\in\mathbb Z[D]$ satisfies $W\equiv D\pmod2$
and $w=W(4)$ is divisible by four.
Write $W=D+2R$. Then $R(4)=(w-4)/2$ is even,
and $R(-2)\equiv R(4)\pmod2$. Therefore
\[
                         W(-2)\equiv2\pmod4.        \tag{G036.U1}
\]
In particular $W$ and $D+2$ are coprime over
$\mathbb Q[D]$.
If $(D+2)G=W F$ for integral polynomials, then
$W\mid G$ in $\mathbb Z[D]$, even when $W$ is
not primitive. To check the content explicitly,
write $W=cW_0$ with $W_0$ primitive and $c>0$.
Gauss's content identity and primitivity of $D+2$
first give $c\mid G$. Write $G=cG_0$.
Coprimality gives $W_0\mid G_0$ over $\mathbb Q[D]$,
and primitivity makes this quotient integral.

Consequently, in any of the closed tables under
consideration, $u_i=6$ forces $U_i=D+2$, and
\[
                 4\mid w,\ u_i=6
                 \quad\Longrightarrow\quad w\mid g_j
                 \qquad(i\ne j).                  \tag{G036.U2}
\]
At $w=12$ this excludes both $(g_1,g_2)=(72,112)$
rows, since $u_1=6$ but $12\nmid112$.
It also excludes the three $(80,96)$ rows with
$d_2=14,u_2=6$, and their three index-reversed rows:
in each case the opposite action has mass eighty.
Thus eight entire base rows, with all their auxiliary
increments, are impossible.

\paragraph{The forward row with action masses eighty and ninety-six.}
Consider
\[
 (w,g_1,g_2,d_1,d_2,u_1,u_2)
                    =(12,80,96,12,4,4,12).         \tag{G036.U3}
\]
Here $U_1=D$ and $\delta_2=\psi$.
If $W(0)\ne0$, the integral quotient argument for
$U_1=D$ gives
$G_2=DWC$ with $C\equiv D-1\pmod2$.
It applies to nonprimitive $W$ as well by the same
content argument just proved. Its value would be
$C(4)=96/(4\cdot12)=2$, whereas its displayed
parity requires an odd value. This is impossible.

It remains that $W=DH$, with the mass-twelve
classification proved above: $H(4)=3$,
$H\equiv1\pmod2$, and $H(0)\ne\pm1$.
Either $H=3$, or $H$ is primitive irreducible
of leading coefficient two.
Since three does not divide $g_1=80$, in the
nonconstant case $H\nmid G_1$.
The second cell identity then forces $H\mid U_2$.
As in the earlier mass-twelve proof, $H(0)$ is
odd and not $\pm1$, whereas $U_2(0)$ is zero
or $\pm8$; therefore $U_2(0)=0$.
Writing $U_2=DJ$, both $H,J$ are primitive
irreducible of leading coefficient two and mass
three, so $J=H$ and $U_2=W$.
For $H=3$, reducing the same identity modulo
three instead gives $3\mid U_2$ coefficientwise;
its effective mass twelve and parity force
$U_2=3D=W$. Thus $U_2=W$ in every case.

The two cell identities now imply
\[
 G_2=G_1+D\psi,\qquad
                    (H-1)G_1=D(H\delta_1+\psi).     \tag{G036.U4}
\]
Write $\delta_1=\psi+2R$, where $R$ is effective
of mass four. Then $\delta_1(0)\in\{-8,0,8\}$.
If $\delta_1(0)=0$, both $\delta_1$ and $\psi$
are divisible by $D$. Since $H(0)\ne1$,
(G036.U4) forces $D^2\mid G_1$.
The quotient has even coefficients but value
$80/16=5$ at four, a contradiction.

If $\delta_1(0)=\pm8$, first $D\mid G_1$.
Divide (G036.U4) by $D$ and evaluate at zero.
Coprimality of the consecutive integers
$H(0),H(0)-1$ shows that $H(0)-1$ divides eight.
The explicit root classification gives
\[
 H(0)=1+2o^2\quad\hbox{or}\quad
 H(0)=1+2(o^2-e^2),
\]
with $o$ positive odd and $e$ positive even.
In particular $H(0)\equiv3\pmod8$.
The only such odd integer with $H(0)-1\mid8$
is three. The first formula then forces $o=1$;
the second would require $o^2-e^2=1$,
impossible for positive $e$.
Consequently $H=3$ and $W=3D$.

Equation (G036.U4) becomes
$G_1=D(2\psi+3R)$.
Its even parity forces $R=2P$ integrally, where
$P$ is effective of mass two, hence either two
or some $T_k$. Thus
\[
 \delta_1=\psi+4P,\qquad
 G_1=D(2\psi+6P),\qquad
 G_2=D(3\psi+6P).
\]
The central Laurent coefficient of $2D\psi$
is four. For $P=T_k$ with $k\ge2$, that of
$6DP$ is zero, whereas $G_1$ must have central
coefficient at least sixteen.
Therefore $P=2$ or $T_1$.

Put $L=2\psi+6P$, an integral quadratic with
leading coefficient two and $L(4)=20$.
The second auxiliary identity gives
$V_2L=3(DG_2-A_2)$.
Modulo three, $L$ reduces to $2\psi\ne0$,
so $V_2=3Z$ integrally.
Since $V_2$ is effective and $V_2\equiv D\pmod2$,
$Z$ is effective and $Z\equiv D\pmod2$.
The actual second auxiliary coefficient is
\[
              A_2=D^2(3\psi+6P)-Z(2\psi+6P).       \tag{G036.U5}
\]
The first term has degree four and leading
coefficient three. A $Z$ of degree greater
than two makes the leading coefficient of
$A_2$ negative. If $\deg Z=2$, its leading
coefficient is positive even and at least two,
so the degree-four coefficient is at most
$3-2\cdot2<0$.
Thus effectivity forces $\deg Z\le1$.
But then $V_2=3Z$ has no $q^2$ coefficient,
contradicting its required coefficient at least two.
The forward row (G036.U3) is excluded.
\hfill$\square$

\paragraph{Geometric scope.}
The divisibility lemma removes ordinary-dimension
patterns for effective graded sign-representation
objects in the specified closed kernel presentation.
It does not assume positive polynomial quotients:
content and integrality alone supply the character
divisibilities. The forward-row exclusion further
uses actual coefficientwise effectivity of the
structure objects. These are the same mathematical
objects of NN-G010 and SGA~3, I, 4.7.2--4.7.3,
with composition in the supplied archive
cote~2, scan~17 transcription. No change of
multiplication-table presentation is permitted
in this conclusion.

\paragraph{The two mass-twelve rows of relaxed cost $2604$.}
Let $G_-$ have mass eighty and $G_+$ have mass ninety-six,
in either orientation of the closed table. The two rows now considered have
\[
 W(4)=12,\qquad \delta_- =\psi+2R,\quad U_-=D+2,
 \qquad \delta_+=\psi,\quad U_+(4)=12,
\]
where $\psi=D(D-3)$ and $R$ is either the constant two or $T_k$,
with $k\ge1$. Every allowed auxiliary increment is included.
In particular $\delta_-(0)=4$ or $-4$.

We first record the divisibility interface used in both orientations.
Write $W=D+2B$ with $B\in\mathbb Z[D]$. Since $W(4)=12$,
$B(4)=4$; therefore $B(-2)$ is even, and
$W(-2)\equiv2\pmod4$. Thus $W$ is coprime to $D+2$
over $\mathbb Q[D]$. From the minus-cell identity
\[
                 (D+2)G_+=W(G_--D\delta_-)
\]
it follows that $W$ divides $G_+$ in $\mathbb Z[D]$.
Here integrality follows either by Gauss's lemma and content,
or directly as follows. Write $W=cW_0$, with $W_0$ primitive.
The rational coprimality gives $G_+=W_0Q$ with $Q$ integral.
After cancellation, $c$ divides every coefficient of $(D+2)Q$;
the product-content identity and primitivity of $D+2$ imply
that $c$ divides every coefficient of $Q$.

If $W(0)\ne0$, the preceding mass-twelve classification gives
$W(0)=\pm8$ and $U_+(0)\in\{-8,0,8\}$.
The two cell identities at zero force both $G_i(0)=0$:
their coefficient determinant is
$W(0)^2-U_-(0)U_+(0)=64-2U_+(0)\ne0$.
Write $G_\pm=DP_\pm$ integrally. Dividing the identities by
$D$ and evaluating at zero yields
\[
 wP_-(0)-2P_+(0)=w\delta_-(0),\qquad
 wP_+(0)-uP_-(0)=0,
 \quad w=W(0),\ u=U_+(0).
\]
Consequently
\[
 P_-(0)=\frac{64\delta_-(0)}{64-2u}.
\]
For $u=8$ or $-8$ this is respectively
$4\delta_-(0)/3$ or $4\delta_-(0)/5$, neither an integer.
Hence $u=0$ and $P_+(0)=0$, so $D^2$ divides $G_+$.
On the other hand $G_+=WQ$ with $Q(4)=96/12=8$.
Since $W(0)\ne0$, $D^2$ divides $Q$ integrally.
Evaluation at four would make sixteen divide eight, a contradiction.

It remains to treat $W(0)=0$, so $W=DH$ with $H$ as in
the preceding mass-twelve gluing lemma. If $H$ is nonconstant,
it is primitive irreducible with leading coefficient two,
$H(4)=3$, and odd $H(0)\ne1,-1$.
It cannot divide $G_-$, since three does not divide eighty.
The plus-cell identity therefore forces $H\mid U_+$.
Divisibility at zero, together with $U_+(0)\in\{-8,0,8\}$,
forces $U_+(0)=0$. Write $U_+=DJ$.
The same classification makes $J$ either three or a primitive
irreducible polynomial of leading coefficient two and value three
at four. Since $H\mid J$, the constant case is impossible and
$J=H$. Thus $U_+=W$.
If $H=3$, reduce the plus-cell identity modulo three.
The reduction of $G_-$ is nonzero because its mass is eighty,
so all coefficients of $U_+$ are divisible by three.
Its odd coefficients at $q,q^{-1}$ are at least three each,
and its even central coefficient at least six. These exhaust
mass twelve, whence again $U_+=3D=W$.

Cancelling $W$ in the plus-cell identity gives
$G_+=G_-+D\psi$. The other cell identity becomes
\[
 (W-D-2)G_-=D\bigl(W\delta_-+(D+2)\psi\bigr).
\]
Its right side is divisible by $D^2$, whereas the first factor
on the left has value $-2$ at zero. Therefore $D^2\mid G_-$.
If $G_-=G_1$, its quotient by $D^2$ has every coefficient even
but evaluation $80/16=5$ at four, a contradiction.
If $G_-=G_2$, the necessary congruence of NN-G030 is
\[
 G_--D\delta_-\equiv G_+\pmod4.
\]
Substituting $G_+=G_-+D\psi$ gives
$2D(\psi+R)\equiv0\pmod4$, so $R\equiv\psi\pmod2$.
This is impossible. The constant two reduces to zero; if
$R=T_k$, its monic degree is $k$, so equality with the degree-two
polynomial $\psi$ would require $k=2$. But
$T_2\equiv D^2$ and $\psi\equiv D^2+D$ modulo two.
Both orientations are excluded for arbitrary grading support.
\hfill$\square$

The objects and scope are those of NN-G036: actual effective graded
sign representations composing in the symmetric closed table.
Restriction at $q=-1$, ordinary dimension at $q=1$, and integral
character divisibility are operations in the torus representation
ring of SGA~3, I, 4.7.2--4.7.3. No quotient is assumed effective,
and no statement about extra-position tables is inferred.

\paragraph{The reverse mass-twelve row with a mass-four action.}
We now exclude the remaining orientation
\[
 (w,g_1,g_2,d_1,d_2,u_1,u_2)
                  =(12,96,80,4,12,12,4).
\]
Write $G_+=G_1$, $G_-=G_2$, $U_+=U_1$, $U_-=U_2=D$,
$\delta_+=\psi=D(D-3)$, and $\delta_-=\psi+2R$.
Here $R$ is an effective symmetric character of mass four.
In particular $|\delta_-(0)|\le8$ and $|U_+(0)|\le8$.

Suppose first $W(0)\ne0$, so $W(0)=\pm8$.
The minus-cell identity
$DG_+=W(G_--D\delta_-)$ and coprimality of $W,D$
give $G_+=WQ$ with $Q$ integral, including by the content
argument if $W$ is nonprimitive.
Then $G_-=D(Q+\delta_-)$.
The plus-cell identity at zero implies $Q(0)=0$, because
its left side is $U_+(0)G_-(0)=0$ and its right side is
$W(0)^2Q(0)$.
Thus $Q=DC$ and
\[
 G_+=DWC,\qquad G_-=D\delta_-+D^2C,\qquad C(4)=2.
\]
The trace parity $G_1\equiv0$, together with $W\equiv D$
modulo two, forces $C\equiv0\pmod2$.
Dividing the plus-cell identity by $D$ now gives
\[
 U_+(\delta_-+DC)=W(WC-\psi).
\]
At zero this says
$64C(0)=U_+(0)\delta_-(0)$.
The right side has absolute value at most sixty-four.
Therefore $|C(0)|\le1$; its even parity forces $C(0)=0$.
The integral polynomial $C$ would be divisible by $D$,
contradicting $C(4)=2$.

Suppose next $W(0)=0$, so $W=DH$ with the previously
proved mass-twelve classification. Since $H(4)=3$ and
$G_-(4)=80$, the same irreducibility/content argument
as in the preceding mass-twelve rows forces $U_+=W$.
More explicitly, a nonconstant $H$ is primitive irreducible
and cannot divide $G_-$; hence it divides $U_+$.
Its odd value $H(0)\ne\pm1$ forces $U_+(0)=0$,
and the classification of $U_+/D$ gives $U_+=W$.
When $H=3$, reduction modulo three and the mass-twelve
effectivity/parity bound give the same conclusion.
Consequently
\[
 G_+=G_-+D\psi,\qquad
                  (H-1)G_-=D(H\delta_-+\psi).
\]
The necessary NN-G030 congruence is
$G_--D\delta_-\equiv G_+\pmod4$.
It forces $R\equiv\psi\pmod2$.
The effective character $R$ has mass four, while
$\psi=T_1+T_2$ already has four odd Laurent coefficients
and mass four. Hence coefficientwise effectivity and
the mass bound force $R=\psi$ exactly. In particular
$\delta_-=3\psi$ and
\[
                    (H-1)G_-=D\psi(3H+1).
\]
Put $J=(H-1)/2\in\mathbb Z[D]$.
Then $J(4)=1$. It is either the constant one (when $H=3$)
or a monic nonconstant polynomial, since the leading
coefficient of a nonconstant $H$ is two.
The preceding identity becomes
\[
                    J(G_--3D\psi)=2D\psi.
\]
Thus the monic $J$ divides $2D^2(D-3)$ in $\mathbb Z[D]$.
Its only possibilities over $\mathbb Q[D]$ are
$D^a(D-3)^b$ with $0\le a\le2$, $0\le b\le1$;
monicity fixes the scalar. Its value one at four
excludes $a>0$. Therefore $J=1$ or $J=D-3$.

If $J=1$, the identities force
\[
 W=3D,\quad G_-=5D\psi,\quad G_+=6D\psi,\quad
 S=10D^3(D-3)^2.
\]
If $J=D-3$, they force
\[
 W=D(2D-5),\quad G_-=D^2(3D-7),\quad
 G_+=2D^2(2D-5),\quad S=2D^3(3D-7).
\]
In each case $S=G_+G_-/W$ is the actual bottom
structure coefficient. Its prescribed coefficientwise
lower bound is $4D^2(D-1)^2$, whose central Laurent
coefficient is $144$.
Using $[q^0]D^h=\binom{2h}{h}$, the two forced
central coefficients of $S$ are respectively
\[
 10(252-6\cdot70+9\cdot20)=120,\qquad
                    6\cdot70-14\cdot20=140.
\]
Both are smaller than $144$, contradicting effectivity
and the given trace domination. This excludes the reverse
row and every permitted auxiliary increment, without a
support-degree restriction.
\hfill$\square$

This is an integral-character obstruction in the same
closed kernel presentation as NN-G036. The final inequality
compares actual multiplicities of the trivial torus character;
the dictionary is SGA~3, I, 4.7.2--4.7.3 and NN-G010.
No conclusion about a different multiplication table follows.

% Final NN-G036 row exclusion and exact conclusion; root whole-proof PASS 2026-11-22T00:21:30Z.
\paragraph{The last mass-twelve row.}
The only row not covered by the preceding arguments is
\[
 w=12,\quad g_1=g_2=96,\quad s=768,\quad
 \delta_1=\delta_2=\psi,\quad u_1=u_2=10.
\]
Every effective $U_i\ge D$ of mass ten with the required
parity is $D+6$ or $D+2+2T_j$. In particular
$U_i(0)\in\{6,-2\}$, so both values are nonzero.

If $W(0)\ne0$, it is $8$ or $-8$ by the mass-twelve
classification. In the pair of cell identities
\[
 WG_1-U_1G_2=DW\psi,\qquad
 WG_2-U_2G_1=DW\psi,
\]
the coefficient determinant at zero is
$64-U_1(0)U_2(0)\ne0$, because the possible products
are $36,-12,4$. Hence both $G_i$ are divisible by $D$.
After cancelling this factor, the right sides are
$W\psi$, still divisible by $D$.
Applying the same invertible matrix at zero again
shows that both $G_i$ are divisible by $D^2$.
Since $WS=G_1G_2$ and $W(0)\ne0$, it follows that
$D^4\mid S$ integrally.
The quotient has even coefficients because $S\equiv0$
modulo two, but its value at four is
$768/4^4=3$. This is impossible.

It remains that $W=DH$, where $H(4)=3$,
$H\equiv1\pmod2$, and $H$ is either three
or primitive irreducible with leading coefficient two.
If $H$ is nonconstant, it cannot divide either
$U_i$, since three does not divide ten.
The common-factor argument of (G036.R2) gives
$H\mid G_1,G_2$ integrally.
If $H=3$, reduce both cell identities modulo three.
The reductions of the $U_i$ are nonzero, since their
masses are ten; the integral domain $\mathbb F_3[D]$
therefore forces the reductions of both $G_i$ to
vanish. Thus $H\mid G_1,G_2$ also in this case.
Write $G_i=HP_i$. Then
\[
 P_i(4)=32,\qquad
 P_1\equiv0,\quad P_2\equiv D^2(D-1)\pmod2,
\]
and the divided cell identities are
\[
 DH P_1-U_1P_2=D^2\psi,\qquad
 DH P_2-U_2P_1=D^2\psi.
\]
Because each $U_i(0)\ne0$, evaluation at zero
forces both $P_i(0)=0$.
Writing $P_i=DQ_i$ and cancelling one $D$ gives
\[
 DH Q_1-U_1Q_2=D\psi,\qquad
 DH Q_2-U_2Q_1=D\psi.
\]
Evaluation at zero again forces both $Q_i(0)=0$.
Consequently $P_i=D^2R_i$ with integral $R_i$.
Their masses are $R_i(4)=32/16=2$.
But $P_2\equiv D^2(D-1)$ forces
$R_2\equiv D-1\pmod2$, whose evaluation at four
is odd. This contradicts $R_2(4)=2$.
The last row is excluded for every grading support.
\hfill$\square$

\paragraph{Exact closed-table optimum and geometric conclusion.}
The complete necessary mass gate of NN-G035,
followed by all the uniform exclusions of NN-G036,
leaves exactly the unique $2744$-pair realization
of NN-G033--NN-G034. Therefore the minimum number
of cancelling trivial/sign pairs in the specified
symmetric closed-table presentation is exactly
$2744$, and its polynomial parameter realization
is unique within that presentation.
All cell degrees and contents, all grading support
locations, and all rational auxiliary parameters
were allowed throughout, subject only to integrality
and effectivity of the actual structure constants.
The final certificate projects every surviving
ordinary-dimension row; it does not enumerate
bounded grading degrees to justify the exclusions.

In the NN-G010 category this is an exact minimum
for actual effective graded sign-representation
objects composing by the fixed kernel table.
The grading is the diagonalizable-group dictionary
of SGA~3, I, 4.7.2--4.7.3, and composition is
the supplied archive cote~2, scan~17 transcription.
The theorem concerns this presentation, not all
possible effective multiplication tables.
The unrestricted extra-position optimization still
has the actual upper bound $2744$, the finite
equal-score cutoff of NN-G029/NN-G033, and the
total-mass FPT algorithm of NN-G032; its numerical
minimum remains open.

